\uuid{fIYn}
\exo7id{1629}
\titre{exo7 1629}
\auteur{barraud}
\organisation{exo7}
\datecreate{2003-09-01}
\isIndication{false}
\isCorrection{false}
\chapitre{Réduction d'endomorphisme, polynôme annulateur}
\sousChapitre{Diagonalisation}
\cpge{mpsi-mp2i.changements-bases.a, pcsi.matrices-applications-lineaires.c, ptsi.matrices-applications-lineaires.c, mpsi-mp2i.determinants.d, pcsi.calcul-matriciel-systemes-lineaires.d, ptsi.calcul-matriciel-systemes-lineaires.d, bcpst1.matrices, mp-mpi.reduction-endomorphismes.b, pc.reduction.a, psi.reduction-endomorphismes.a, pt.reduction-endomorphismes.a, mp-mpi.reduction-endomorphismes.c, pc.reduction.b, psi.reduction-endomorphismes.b, pt.reduction-endomorphismes.b, mp-mpi.reduction-endomorphismes.d, pc.reduction.c, psi.reduction-endomorphismes.c, pt.reduction-endomorphismes.c, mp-mpi.reduction-endomorphismes.e, pc.reduction.e, psi.reduction-endomorphismes.e, pt.reduction-endomorphismes.d, bcpst2.valeurs-propres-vecteurs-propres.a, bcpst2.valeurs-propres-vecteurs-propres.b}
\module{Algèbre}
\niveau{L2}
\difficulte{}

\contenu{
\question{
Les matrices suivantes sont-elles diagonalisables,
triangularisables ? Si oui, les réduire.
$$
A_{1}=\begin{pmatrix}
        3 &-1 & 1 \\
        2 & 0 & 1 \\
        1 &-1 & 2
      \end{pmatrix}
\qquad
A_{2}=\begin{pmatrix}
        3 & 2 &-2 \\
       -1 & 0 & 1 \\
        1 & 1 & 0
      \end{pmatrix}
\qquad
A_{3}=\begin{pmatrix}
       13 &-5 &-2 \\
       -2 & 7 &-8 \\
       -5 & 4 & 7
      \end{pmatrix}
\qquad
$$
}
}
